%% kapa: exam=2 round=30 year=2019 period=3 subject=compensation_law subject_ko=감정평가 및 보상법규 seq=2 number=2 points=30
%% source: Q-Net 감정평가사 기출문제 (cst00309) artlSeq=5205595 file=raw/1008/5205595/감정평가 및 보상법규.pdf
%% conversion: kordoc 4.24.51 -> tools/gen2.py (auto); review=proofread-gpt-terra; flags=none
\begin{longquestion}{【문제 2】}
甲은 골프장을 보유ㆍ운영해왔는데, 그 전체 부지 1,000,000 m\textsuperscript{2} 중 100,000 m\textsuperscript{2}가 도로건설 사업부지로 편입되었고, 골프장은 계속 운영되고 있다. 위 사업부지로 편입된 부지 위에는 오수처리시설이 있었는데, 수용재결에서는 그 이전에 필요한 비용으로 1억 원의 보상금을 산정하였다. 다음 물음에 답하시오. (30점) \par
\medskip\noindent 물음 1) 甲은 골프장 잔여시설이 종전과 동일하게 운영되려면 위 오수처리시설을 대체하는 새로운 시설의 설치가 필요하다고 보아 그 설치에 드는 비용 1억 5천만 원을 보상받아야 한다고 주장한다. 甲의 주장은 법적으로 타당한가? (10점) \par
\medskip\noindent 물음 2) 甲은 골프장 잔여시설의 지가 및 건물가격 하락분에 대하여 보상을 청구하려고 한다. 이때 甲이 제기할 수 있는 소송에 관하여 설명하시오. (20점) \par
\begin{framedtext}
참조 조문 \par
< 공익사업을 위한 토지 등의 취득 및 보상에 관한 법률 > \par
제73조(잔여지의 손실과 공사비 보상) ① 사업시행자는 동일한 소유자에게 속하는 일단의 토지의 일부가 취득되거나 사용됨으로 인하여 잔여지의 가격이 감소하거나 그 밖의 손실이 있을 때 또는 잔여지에 통로ㆍ도랑ㆍ담장 등의 신설이나 그 밖의 공사가 필요할 때에는 국토교통부령으로 정하는 바에 따라 그 손실이나 공사의 비용을 보상하여야 한다. 다만, 잔여지의 가격 감소분과 잔여지에 대한 공사의 비용을 합한 금액이 잔여지의 가격보다 큰 경우에는 사업시행자는 그 잔여지를 매수할 수 있다. \par
②～⑤ 생략 \par
제75조의2(잔여 건축물의 손실에 대한 보상 등) ① 사업시행자는 동일한 소유자에게 속하는 일단의 건축물의 일부가 취득되거나 사용됨으로 인하여 잔여 건축물의 가격이 감소하거나 그 밖의 손실이 있을 때에는 \par
국토교통부령으로 정하는 바에 따라 그 손실을 보상하여야 한다. 다만, 잔여 건축물의 가격 감소분과 보수비(건축물의 나머지 부분을 종래의 목적대로 사용할 수 있도록 그 유용성을 동일하게 유지하는 데에 일반적으로 필요하다고 볼 수 있는 공사에 사용되는 비용을 말한다. 다만, 「건축법」 등 관계 법령에 따라 요구되는 시설 개선에 필요한 비용은 포함하지 아니한다)를 합한 금액이 잔여 건축물의 가격보다 큰 경우에는 사업시행자는 그 잔여 건축물을 매수할 수 있다. \par
②～⑤ 생략 \par
\end{framedtext}
\end{longquestion}
